\section{Inputs and the radial conversion}
We use the notation of the construction above. To state the consequences on both
spaces together, let
\[
(\mathcal X,M,k,P_k)=
\begin{cases}
(c_0(I),A,g,P),&I=\{0,1,\ldots\}\times\N,\\
(\ell^1(\N),T,f,B),&f=Q^*g,
\end{cases}
\qquad P_kx=\pair{x}{k}k.
\]
The duals are respectively the full $\ell^1(I)$ and $\ell^\infty(\N)$.
We retain the gap and split conventions
\[
\Phi_M(z,p)=\sup_{(x,a)\in\gra M}\pair{z-x}{a-p},\qquad
\Delta_{M,S}(z,p)=\inf_{b\in\mathcal X^*}
       [\Phi_M(z,p-b)+\Phi_S(z,b)].
\]
Theorems~\ref{thm:C} and~\ref{thm:L} and
Eqs.~\eqref{eq:C2}, \eqref{eq:C9}, and~\eqref{eq:L6} prove
that $M$ is maximally monotone and that every $(x,a)\in\gra M$ satisfies
\begin{equation}
\begin{gathered}
t=\pair{x}{k},\qquad |t|>1,\qquad
\|k\|=2,\qquad \|x\|>\tfrac12,\\
\pair{x}{a}=t^2-1-W(|t|-1),\qquad
0\le W(s)\le\sigma=\frac{\pi^2}{96}<\frac18.
\end{gathered}
\label{eq:cor-input}
\end{equation}
Here $W(s)=\sum_{n\ge1}(\max\{1-sn^{1/4},0\}/(4n))^2$ for $s>0$.
Every $|t|>1$ is realized by an actual graph point: on $c_0(I)$ use the
label $a^t$ in Eq. (D9); on $\ell^1$ use an actual lift through the
surjection $Q$. No bounded linear right inverse is used. The preceding construction
also proves maximality of $P_k$ and the exact split value
\begin{equation}
\Delta_{M,P_k}(0,0)=\sigma.
\label{eq:cor-split-input}
\end{equation}
Only the final finite-split assertion in Corollary~\ref{cor:universal}
uses Eq.~\eqref{eq:cor-split-input}.

For any maximal $M$ on a real Banach space, write
\begin{equation}
L_M(z,p)=\max\left\{0,
\sup_{\substack{(x,a)\in\gra M\\x\ne z}}
\frac{\pair{p-a}{x-z}}{\|x-z\|}\right\}.
\label{eq:cor-Verona}
\end{equation}
We use the nonnegative convention of \cite{verona2008}. The displayed
definition in \cite[Section~2]{verona2012} omits the outer maximum with
zero. At a missing zero primal, the notation in
the preceding construction satisfies exactly
\begin{equation}
\mathcal V(\gra M)=
\sup_{(x,a)\in\gra M}\frac{\max\{0,-\pair{x}{a}\}}{\|x\|}
=L_M(0,0).
\label{eq:cor-V}
\end{equation}
The definition is not a finiteness theorem. For the two concrete
operators, Eq.~\eqref{eq:cor-input} gives, on every graph row,
\[
\pair{x}{a}>-\sigma,\qquad
\frac{\max\{0,-\pair{x}{a}\}}{\|x\|}
\le\frac{\sigma}{\|x\|}<2\sigma.
\]
Taking the supremum proves $\mathcal V(\gra M)\le2\sigma<1/4$.
This bound comes from the construction, not from a universal finiteness
claim or an assumed failure of the sum conjecture.

\begin{lemma}[Radial conversion]\label{lem:radial-conversion}
Let $M:\mathcal X\rightrightarrows\mathcal X^*$ be maximally monotone
on a real Banach space, with $0\notin\dom M$. For $\lambda>0$, put
$J_\lambda=\partial(\lambda\|\cdot\|)$. Then
\begin{equation}
(0,0)\in\gra(M+J_\lambda)^\mu\setminus\gra(M+J_\lambda)
\quad\Longleftrightarrow\quad
\mathcal V(\gra M)\le\lambda.
\label{eq:cor-exact-radial}
\end{equation}
Whenever this holds, $(M,J_\lambda)$ satisfies the original intersection
condition and its sum is not maximally monotone.
\end{lemma}
\begin{proof}
The norm subdifferential has the formula
\[
J_\lambda x=\{j\in\mathcal X^*: \|j\|\le\lambda,
                         \pair{x}{j}=\lambda\|x\|\}.
\]
Hahn--Banach makes this set nonempty for every $x$; at zero it is the
closed dual ball of radius $\lambda$. It is maximally monotone by the
single-subdifferential theorem \cite[Theorem~A]{rockafellar1970subdiff}.
For every $(x,a)\in\gra M$ and every $j\in J_\lambda x$,
\[
\pair{x}{a+j}=\pair{x}{a}+\lambda\|x\|.
\]
Thus relatedness to $(0,0)$ is equivalent to the right-hand side being
nonnegative on every row, which is exactly
$\mathcal V(\gra M)\le\lambda$. The target is outside the sum graph
because $M(0)=\varnothing$. The sum is monotone, so adjoining that target
gives a proper monotone extension. Finally, $\dom J_\lambda=\mathcal X$
and a maximal monotone graph is nonempty, giving the intersection condition.
\end{proof}
The mechanism in this lemma is the missing-primal specialization of the
radial argument in \cite[Theorem~1.2, proof]{verona2008}; the displayed
proof records the full graph quantifiers and the counterexample checks.

\section{Corollaries and their witnesses}
\begin{corollary}[Norm-subdifferential partners]\label{cor:radial}
For each of the two concrete operators $M$, and every $\lambda\ge2\sigma$,
the pair $(M,J_\lambda)$ consists of maximally monotone operators,
$\dom J_\lambda=\mathcal X$, and $M+J_\lambda$ is not maximally monotone.
The target is $(0,0)$. In particular, $\lambda=1/4$ is permitted.
\end{corollary}
\begin{proof}
Use Eq.~\eqref{eq:cor-V} and Lemma~\ref{lem:radial-conversion}.
The finite bound applies to the whole graph of each operator, including
all kernel translates in the $\ell^1$ construction.
\end{proof}

\begin{corollary}[Failure of the regularity conditions]\label{cor:regularity}
For each of the two concrete operators, all five of the following
conditions fail:
\begin{enumerate}
\item $L_M(z,p)=\dist(p,Mz)$ for every $(z,p)\in\mathcal X\times\mathcal X^*$;
\item $M+S$ is maximally monotone for every bounded-range maximally
monotone $S:\mathcal X\rightrightarrows\mathcal X^*$;
\item $M+\partial(\lambda\|\cdot-z\|)$ is maximally monotone for every
$z\in\mathcal X$ and every $\lambda>0$;
\item $L_M(z,0)<\infty$ implies $z\in\dom M$ for every $z\in\mathcal X$;
\item $L_M(z,0)=\dist(0,Mz)$ for every $z\in\mathcal X$.
\end{enumerate}
These are the nonnegative-convention formulation of the regularity
conditions in \cite[Theorem~2.3]{verona2012}. Their displayed definition
uses a bare supremum; the two conventions agree at our missing-zero
witness, since maximality forces a graph row with negative self-pairing.
We prove each failure directly below. In particular, both operators are
nonregular in the convention of \cite{verona2008}.
\end{corollary}
\begin{proof}
We use $\dist(p,\varnothing)=+\infty$. The first, fourth, and fifth
conditions fail at $z=0$, $p=0$, since
\[
L_M(0,0)\le2\sigma<\infty=\dist(0,M0),\qquad 0\notin\dom M.
\]
For the second and third take $S=J_{1/4}$ and $z=0$ in
Corollary~\ref{cor:radial}. Its range is contained in the closed dual
ball of radius $1/4$. This proves each failure directly. Bounded range
is not being inferred from boundedness of the rank-one linear map $P_k$.
\end{proof}

\begin{corollary}[An explicit finite-gap FPV failure]\label{cor:fpv}
For each concrete $M$, set
\begin{equation}
U=\{x\in\mathcal X:\pair{x}{k}>-1\},\qquad
d=0,\qquad p=-\sigma k.
\label{eq:cor-FPV-witness}
\end{equation}
Then $U$ is open and convex, $0\in U$, $U\cap\dom M\ne\varnothing$,
and $(0,p)\notin\gra M$, whereas
\begin{equation}
\pair{-x}{p-a}>0
\quad\text{for every }(x,a)\in\gra M\text{ with }x\in U,
\qquad
\Phi_M(0,p)=2\sigma>0.
\label{eq:cor-FPV-checks}
\end{equation}
Consequently, each of the same operators fails type (FPV) and its
restriction to points at which the Fitzpatrick gap is finite.
\end{corollary}
\begin{proof}
Continuity of $k$ proves openness and convexity. An actual row with
$t=2$ proves $U\cap\dom M\ne\varnothing$. For any graph row in $U$,
$t>-1$ and $|t|>1$ force $t>1$. Therefore
\begin{align*}
\pair{-x}{p-a}
&=\pair{x}{a-p}=t^2-1-W(t-1)+\sigma t\\
&\ge t^2-1+\sigma(t-1)>0.
\end{align*}
On the entire graph, including the negative branch, put $r=|t|>1$.
Then
\begin{align*}
\pair{-x}{a-p}
&=1+W(r-1)-t^2-\sigma t\\
&\le1+\sigma-r^2+\sigma r\\
&=2\sigma-(r-1)(r+1-\sigma)<2\sigma.
\end{align*}
This proves the upper bound on the supremum. Along actual parameters
$t=-1-s$, $s\downarrow0$, we have $W(s)\to\sigma$, and the displayed
expression tends to $2\sigma>0$. Dominated convergence applies since
each squared coordinate is bounded by $1/(16n^2)$.
Finally, $M(0)=\varnothing$ gives graph nonmembership. Type (FPV)
requires membership whenever an open convex set meets the domain and
the target is locally monotonically related there. The finite-gap
restriction imposes in addition $\Phi_M(d,p)<\infty$. All these premises
have just been checked, so both properties fail.
\end{proof}
This witness meets the actual domain. It is not a vacuous test in a
domain-free ball, and no replacement of the first operator by a product
operator has been made.

\begin{corollary}[A bounded-domain normal-cone partner]\label{cor:normal}
For either concrete $M$, choose an actual primal $x_2\in\dom M$ with
$\pair{x_2}{k}=2$. Fix $\rho>\|x_2\|$ and define
\begin{equation}
C_\rho=\{x\in\mathcal X:\|x\|\le\rho,\ \pair{x}{k}\ge-1/2\}.
\label{eq:cor-cut}
\end{equation}
Then $M$ and $N_{C_\rho}$ are maximally monotone,
$\dom M\cap\operatorname{int}\dom N_{C_\rho}\ne\varnothing$, and
their common domain is bounded. Their sum is not maximally monotone;
$(0,-\sigma k)$ belongs to the monotone polar of the sum graph
and lies outside that graph.
In both spaces one can take $\|x_2\|=8$ and $\rho=9$.
\end{corollary}
\begin{proof}
The actual choices from the core proofs are
\[
x_2=-6e_{0,1}-8e_{0,2}-3e_{0,3}\quad\text{on }c_0(I),
\qquad x_2=8e_{n_0}\quad\text{on }\ell^1,
\]
where $q_{n_0}=-(3/4)e_{0,1}-e_{0,2}-(3/8)e_{0,3}$. Both have norm
$8$, and their pairing with $k$ is $2$. The set $C_\rho$ is closed,
convex, bounded, and contains both $0$ and
$x_2$ in its interior. The normal cone is
\[
N_{C_\rho}(x)=\{n\in\mathcal X^*:
\pair{y-x}{n}\le0\ \text{for all }y\in C_\rho\}
\quad(x\in C_\rho),
\]
and is empty off $C_\rho$. It is the subdifferential of the proper
lower semicontinuous convex indicator, hence is maximally monotone
\cite[Theorem~A]{rockafellar1970subdiff}. Since zero is a normal at
every point of $C_\rho$, its domain is exactly $C_\rho$.
Thus $x_2$ verifies the intersection condition, and the common domain
is contained in the ball of radius $\rho$.

Let $p=-\sigma k$. For every $(x,a)\in\gra M$ with $x\in C_\rho$,
$t\ge-1/2$ and $|t|>1$ imply $t>1$. For every $n\in N_{C_\rho}(x)$,
testing the normal inequality at $y=0$ gives $\pair{x}{n}\ge0$.
The calculation in Corollary~\ref{cor:fpv} now gives
\[
\pair{-x}{p-a-n}=\pair{x}{a-p}+\pair{x}{n}>0.
\]
It holds for every sum value. The target is not in the sum graph because
$0\notin\dom M$, so it gives a proper monotone extension.
\end{proof}
This supplies a concrete bounded-common-domain failure, the restricted
sum assertion shown equivalent to the unrestricted one in
\cite[Theorem~2]{botbuenosimons2016}. No maximality of an intermediate
sum is assumed, and boundedness is not replaced by compactness.

\begin{corollary}[Arbitrarily small positive slope bounds]\label{cor:flat}
For either space and each $\theta>0$, there is a maximally monotone
$M_\theta$ with $0\notin\dom M_\theta$ such that
\[
0<\mathcal V(\gra M_\theta)\le\theta.
\]
For each prescribed $\lambda>0$, one may choose such an operator so
that $M_\theta+J_\lambda$ is a nonmaximal sum satisfying the original
intersection condition. The zero endpoint is impossible for a maximal
operator with missing zero primal.
\end{corollary}
\begin{proof}
For $s>0$, define $M_s(x)=sM(x)$. Dividing a polar query's dual
coordinate by $s$ proves that $M_s$ is maximally monotone; its domain
is unchanged. Direct homogeneity gives
$\mathcal V(\gra M_s)=s\mathcal V(\gra M)$.
Taking $s=\theta/(2\sigma)$ proves the upper bound. A zero value would
make $\pair{x}{a}\ge0$ on the entire graph, so $(0,0)$ would be in
its polar; maximality would put it in the graph, contradicting the
missing zero primal. For a prescribed $\lambda$, take
$\theta=\lambda/2$ and apply Lemma~\ref{lem:radial-conversion}.
\end{proof}
Here the quantifier is every positive bound with a possibly different
scaled operator. It does not assert that one fixed operator has slope
below every positive number. In particular, one may take $\theta=1/2$.

\begin{corollary}[Universal assertions and their negations]\label{cor:universal}
Conditional on the two core theorems, the following universal assertions
are false, already on each of the two displayed spaces:
\begin{enumerate}
\item Every qualified pair of maximally monotone operators has a maximally
monotone sum, including the restriction to bounded common domain.
\item Every maximally monotone operator is regular, or satisfies any of
the five equivalent conditions in Corollary~\ref{cor:regularity}.
\item Every maximally monotone operator has the local graph-membership
property defining type (FPV), even if that implication is required only
at points with finite $\Phi_M(d,p)$, with or without the requirement
$U\cap\dom M\ne\varnothing$ in its premises.
\item Every maximally monotone operator omitting zero has
$\mathcal V(\gra M)=+\infty$.
\item For any prescribed $\lambda>0$, all maximally monotone operators
have a maximally monotone sum with $J_\lambda$.
\item There is no finite-$\Delta$ counterexample under the original
intersection condition.
\end{enumerate}
The corresponding assertions over all separable real Banach spaces,
or over all real Banach spaces, are therefore false as well.
\end{corollary}
\begin{proof}
The witnesses are, respectively, Corollary~\ref{cor:normal},
Corollary~\ref{cor:regularity}, Corollary~\ref{cor:fpv},
Eq.~\eqref{eq:cor-V} with the finite bound derived after it, and
Corollary~\ref{cor:flat}. For the last assertion, use the original
rank-one pairs and Eq.~\eqref{eq:cor-split-input}. Both spaces are
separable, so the same witnesses disprove the larger-carrier assertions.
\end{proof}

\noindent\textit{Scope of equivalence and attribution.}
Negating $\forall M\,\mathcal P(M)$ produces $\exists M\,\neg\mathcal P(M)$,
not $\forall M\,\neg\mathcal P(M)$. A valid equivalence between universal
assertions remains valid when those assertions are false. The five
fixed-operator equivalences cited above are classical. Each failure has
the explicit witness given above. These corollaries use no additional
equivalence between assertions on different operators or spaces.

\section{Consequences for quotient norms and graph representations}
The following consequences use the actual maps of the construction. They
identify the information retained by the dual quotient and the information
lost by extending a curve in the Hilbert comparison space. They are separate deductions,
not additional assumptions in the core maximality proof.

\begin{corollary}[Constant pairing and an unbounded norm in $\ell^\infty/c_0$]\label{cor:corona}
For the standard $\ell^1$ operator $T$ in the construction above, there are
actual rows $(u_N,p_N)\in\gra T$ and a fixed $d\in\ell^1$ such that
\begin{equation}
\|u_N\|_1=8,\qquad \pair{u_N}{p_N}=3,\qquad
\dist_{\infty}(p_N,c_0)=\|p_N\|_\infty=5+2N,
\label{eq:cor-constant-work}
\end{equation}
and
\begin{equation}
\|d\|_1=\frac{\pi^2}{6},\qquad \pair{d}{f}=0,\qquad
\pair{d}{p_N}=-2\sum_{r\ge1}\frac{\min\{N,r\}}{r^2}\longrightarrow-\infty.
\label{eq:cor-fixed-detector}
\end{equation}
The fibre of $(\gra T)^\mu$ above $d$ is empty. In particular, no finite
constants $K,\eta\ge0$ satisfy
\begin{equation}
\dist_{\infty}(p,c_0)\le K+\eta\pair{u}{p}
\quad\text{for every }(u,p)\in\gra T.
\label{eq:cor-WCC-failure}
\end{equation}
\end{corollary}
\begin{proof}
The quotient norm is $\|p+c_0\|_{\ell^\infty/c_0}=\dist_\infty(p,c_0)$.
We first compute it for the enumeration defining $Q$.
For every $a\in\ell^1(I)$,
\begin{equation}
\dist_{\infty}(Q^*a,c_0)
=\limsup_j|\pair{q_j}{a}|=\|a\|_1.
\label{eq:cor-quotient-isometry}
\end{equation}
Indeed, the upper bound follows from $\|q_j\|_\infty\le1$. For
$a\ne0$ and $0<\varepsilon<\|a\|_1$, a finitely supported rational
vector in the unit ball gives pairing greater than
$\|a\|_1-\varepsilon$ in absolute value. Rational scalar perturbations
tending to one supply infinitely many distinct such vectors (start
with a strict bound). Every one occurs in the enumeration, so the
same lower bound holds after deleting any finite prefix. Let
$\varepsilon\downarrow0$. The case $a=0$ is immediate. For any bounded
sequence, its distance to $c_0$ equals the limsup of the absolute
coordinates: truncation gives one inequality, and subtracting any
norm-null sequence gives the other. This proves
Eq.~\eqref{eq:cor-quotient-isometry}, including
$\|Q^*a\|_\infty=\|a\|_1$.

Use actual labels
\[
a_N=-2e_{0,1}+3e_{0,3}
       +\sum_{b=1}^N(e_{b,1}-e_{b,2}).
\]
They have $r(a_N)=2$ and $ma_N=(1,0,\ldots)=F(2)$, since
$\omega_1=0$. Thus $a_N\in D$. Directly from the triangular map,
\[
x_N=La_N=-6e_{0,1}-8e_{0,2}-3e_{0,3}
                  +\sum_{b=1}^N(e_{b,1}+e_{b,2}).
\]
Consequently $\|x_N\|_\infty=8$, $\pair{x_N}{a_N}=3$, and
$\|a_N\|_1=5+2N$. Let $i_N$ be the first index with
$q_{i_N}=x_N/8$, and put $u_N=8e_{i_N}$ and $p_N=Q^*a_N$.
These are actual rows because $Qu_N=x_N$. The pairing is preserved,
and Eq.~\eqref{eq:cor-quotient-isometry} gives
Eq.~\eqref{eq:cor-constant-work}.

For $r\ge1$, the finite primal vector
$v_r=\sum_{b=1}^r(e_{b,1}-e_{b,2})$ is rational and has norm one.
Let $j_r$ be its first index in the enumeration, so the $j_r$ are
distinct. Set $d=-\sum_{r\ge1}r^{-2}e_{j_r}$. This converges in
$\ell^1$, has norm $\pi^2/6$, and
\[
f(j_r)=\pair{v_r}{g}=0,\qquad
p_N(j_r)=\pair{v_r}{a_N}=2\min\{N,r\}.
\]
For each fixed $N$ the series defining $\pair{d}{p_N}$ converges
absolutely. Its negative is at least
$2\sum_{r=1}^N1/r$, proving Eq.~\eqref{eq:cor-fixed-detector}.
Every primal in $\dom T$ has $|\pair{u}{f}|>1$, whereas
$\pair{d}{f}=0$. Thus $d\notin\dom T$, and maximality
$(\gra T)^\mu=\gra T$ makes its entire polar fibre empty.
Finally, Eq.~\eqref{eq:cor-constant-work} would turn
Eq.~\eqref{eq:cor-WCC-failure} into $5+2N\le K+3\eta$ for every
$N$, which is impossible.
\end{proof}
The fixed vector $d$ detects divergence of the dual sequence, although
all selected graph points have the same primal norm and pairing.
Since $d\notin\dom T$, no point of $(\gra T)^\mu$ has primal coordinate
$d$. In particular, $\dist_\infty(p,c_0)$ cannot be bounded on $\gra T$
by a function of $\|u\|_1$ and $\pair{u}{p}$ that is finite at $(8,3)$.

\begin{corollary}[A Lipschitz extension with no additional preimages in the dual space]\label{cor:dense-profile}
For the displayed $c_0(I)$ construction, the map $(r,m)$ has range
\begin{equation}
\operatorname{range}(r,m)=\R\times\ell^1(\{0,1,\ldots\})
\subsetneq\R\times\ell^2(\{0,1,\ldots\}),
\label{eq:cor-dense-range}
\end{equation}
which is dense. The profile $F:J\to\ell^2$ has a unique global
$1$-Lipschitz extension, namely
\begin{equation}
\widehat F(t)=
\begin{cases}
F(t),&|t|>1,\\
(t,\omega_0),&|t|\le1.
\end{cases}
\label{eq:cor-profile-extension}
\end{equation}
Its graph strictly extends $\gra F$, and is maximal for the quadratic
positivity inequality $\|y-y'\|_2^2\le|t-t'|^2$ on
$\R\times\ell^2$. Nevertheless, its entire actual pullback satisfies
\begin{equation}
\{(La,a):a\in\ell^1(I),\ ma=\widehat F(r(a))\}
=\gra A.
\label{eq:cor-same-pullback}
\end{equation}
For these maps, a maximal monotone graph can therefore have a nonmaximal
image under the quadratic representation described in the proof, despite
the dense range of $(r,m)$.
\end{corollary}
\begin{proof}
Every $ma$ is summable. Conversely, for any $t\in\R$ and any
$y\in\ell^1(\{0,1,\ldots\})$, the actual label
\[
a=-t e_{0,1}+(t+y_0)e_{0,3}+\sum_{n\ge1}y_n e_{n,1}
\]
satisfies $r(a)=t$ and $ma=y$. This proves
Eq.~\eqref{eq:cor-dense-range}; strictness follows from
$\omega_0\in\ell^2\setminus\ell^1$, and density from finite truncations.

The coordinatewise positive-part estimate gives
\[
\|\omega_s-\omega_q\|_2^2\le c|s-q|^2,\qquad
c=\frac1{16}\sum_{n\ge1}n^{-3/2}<1,
\]
also for $s=0$ or $q=0$. Put $\chi(t)=\max\{-1,\min\{1,t\}\}$ and
$\zeta(t)=\max\{|t|-1,0\}$. Then
$\widehat F(t)=(\chi(t),\omega_{\zeta(t)})$.
The same-branch and opposite-branch cases give
\[
|\chi(t)-\chi(u)|+|\zeta(t)-\zeta(u)|\le|t-u|.
\]
Therefore
\[
\|\widehat F(t)-\widehat F(u)\|_2^2
\le|\chi(t)-\chi(u)|^2+c|\zeta(t)-\zeta(u)|^2
\le|t-u|^2.
\]
Any global $1$-Lipschitz extension must take the values
$(-1,\omega_0)$ and $(1,\omega_0)$ at the two endpoints, by
continuity from the actual branches. Their distance is $2$.
For an intermediate $t$, the bounds $t+1$ and $1-t$ on the two
distances add to $2$. Equality in the Hilbert-space triangle inequality
forces the affine segment value $(t,\omega_0)$. Hence the extension
is unique. Its graph is maximal for the displayed quadratic inequality:
testing any related $(t,y)$ against $(t,\widehat F(t))$ forces
$y=\widehat F(t)$.

No new parameter $|t|\le1$ has an actual label, because
$(t,\omega_0)\notin\ell^1$. Thus the label constraint for
$\widehat F$ is exactly the original constraint for $F$, proving
Eq.~\eqref{eq:cor-same-pullback}.

We make the last correspondence assertion precise. Set $V=\R$,
$H=\ell^2(\{0,1,\ldots\})$, and $\mathcal T v=vh$. The core identities are
\begin{align*}
\pair{La}{b}+\pair{Lb}{a}&=2[r(a)r(b)-(ma,mb)_H],\\
\pair{\mathcal T v}{a}&=v r(a),\qquad
(\ker r\cap\ker m)^\perp\cap c_0(I)=\R h.
\end{align*}
For the last identity, annihilating differences within every positive
block makes each block constant; the $c_0$ condition makes it zero.
In block zero the tail similarly vanishes and the first two entries
are equal. The converse follows from $r(k)=0$ on the kernel.
Thus the bilinear and adjoint identities and the displayed annihilator
identity all hold.
On $x=Lp+2\delta h$ the coordinate map is
\[
J_{\rm coord}(x,p)=(r(p)+\delta,mp,\delta),\qquad
q(t,y,\delta)=t^2-\|y\|_2^2-\delta^2.
\]
For this quadratic form, a nonempty set $C$ is $q$-positive if
$q(c-c')\ge0$ for all $c,c'\in C$. Define
$C^q=\{z:q(z-c)\ge0\text{ for every }c\in C\}$.
Maximality is with respect to inclusion among $q$-positive sets.
The image of the actual maximal graph is
$C^0=\{(t,F(t),0):t\in J\}$. It is not maximally $q$-positive,
since $\{(t,\widehat F(t),0):t\in\R\}$ properly extends it and
is $q$-positive. In fact, the entire quadratic polar is exactly
\[
(C^0)^q=\{(t,\widehat F(t),0):t\in\R\}.
\]
To verify the reverse inclusion, let $(\tau,y,\delta)$ be related
to $C^0$. If $|\tau|>1$, testing at the same parameter forces
$y=F(\tau)$ and $\delta=0$. If $|\tau|\le1$, limits along the
actual branches give
\[
\|y-(1,\omega_0)\|^2+\delta^2\le(1-\tau)^2,\qquad
\|y-(-1,\omega_0)\|^2+\delta^2\le(1+\tau)^2.
\]
Multiplication by $(1+\tau)/2$ and $(1-\tau)/2$, respectively,
and addition cancel the common term $1-\tau^2$, leaving
$\|y-(\tau,\omega_0)\|^2+\delta^2\le0$.
This also proves uniqueness of the maximal reduced extension, even
when the coordinate $\delta$ is allowed. Nevertheless its new points
have no actual preimages under $J_{\rm coord}$.
Consequently, dense joint range cannot replace actual
joint surjectivity in a correspondence sending every actual maximal
extension to a maximal reduced positive set. This is a failure for
the displayed maps, not an assertion excluding every alternative
representation of $A$.
\end{proof}
The Lipschitz extension is unique. Its values at $|t|\le1$ have no
preimages under $(r,m)$, so the graph defined by the displayed
constraint remains unchanged.

\begin{corollary}[A gap between optimization over the domain and its convex hull]\label{cor:two-row}
On each of the two original spaces there is an operator $\widetilde M$
obtained from $M$ by an affine change of variables, a norm-one functional $\varphi$, and a
unit vector $e$ with $\pair{e}{\varphi}=1$ such that
\begin{equation}
0\le \sup_{\substack{z\in\dom\widetilde M\\\|z\|\le1}}
\pair{z}{\varphi}\le\frac38,
\qquad
\inf_{p\in X^*}\left\{
\sup_{z\in\dom\widetilde M}\pair{z}{p}+\|\varphi-p\|
\right\}=1.
\label{eq:cor-PCT-gap}
\end{equation}
There are two graph points whose mean primal coordinate is $3e/4$.
Its value under $\varphi$ is $3/4$, strictly above the left supremum in
Eq.~\eqref{eq:cor-PCT-gap}. Moreover,
\begin{equation}
\dist(3e/4,\dom\widetilde M)\ge\frac38,
\qquad \Phi_{\widetilde M}(3e/4,0)=\sigma.
\label{eq:cor-finite-projection-hole}
\end{equation}
The transformation changes the operator, but not the Banach space or
its norm.
\end{corollary}
\begin{proof}
Let $x_2$ denote the actual primal of norm $8$ used above, with
$\pair{x_2}{k}=2$. Both $x_2$ and $-x_2$ are actual primals.
Set $\varphi=k/2$. On $c_0(I)$ take
$e=e_{0,1}-e_{0,2}$, regarded as a primal vector. On standard
$\ell^1$ take $e=e_j$ at an index where
$q_j=e_{0,1}-e_{0,2}$. In both cases
$\|e\|=\|\varphi\|=1$ and $\pair{e}{\varphi}=1$.
Put $r_0=3/4$ and define
\[
\mathcal S z=-\frac{x_2}{r_0}\pair{z}{\varphi}
                 +z-\pair{z}{\varphi}e,
\qquad
\widetilde M(z)=\mathcal S^*M(x_2+\mathcal S z).
\]
Here $\mathcal S$ is a bounded linear isomorphism, and the dual variable
is transformed by $\mathcal S^*$. Indeed, writing $z=\alpha e+w$ with $w\in\ker\varphi$
gives $\mathcal S z=-\alpha x_2/r_0+w$ and
$\pair{\mathcal S z}{\varphi}=-\alpha/r_0$. For a prescribed
image $v$, the unique inverse is obtained with
$\alpha=-r_0\pair{v}{\varphi}$ and
$w=v+\alpha x_2/r_0\in\ker\varphi$; it is bounded.
The identity $\pair{z-z'}{\mathcal S^*(a-a')}
=\pair{\mathcal S(z-z')}{a-a'}$ preserves every monotonicity bracket.
The resulting bijection on primal--dual pairs therefore preserves
monotonicity and full maximality.

Every $z\in\dom\widetilde M$ satisfies
\[
1<|\pair{x_2+\mathcal S z}{k}|
  =|2-2\pair{z}{\varphi}/r_0|.
\]
Consequently
\begin{equation}
\pair{z}{\varphi}<\frac38
\quad\text{or}\quad
\pair{z}{\varphi}>\frac98.
\label{eq:cor-transformed-strip}
\end{equation}
The second alternative is impossible when $\|z\|\le1$.
The primals $0$ and $2r_0e=3e/2$ are actual, since their images
are $x_2$ and $-x_2$. This proves the singleton bounds in
Eq.~\eqref{eq:cor-PCT-gap} and provides strict access inside the
unit ball at $0$.

Writing $D'=\dom\widetilde M$, the actual endpoints imply
$e\in\operatorname{conv}D'$. For every $p\in X^*$,
\[
\sup_{z\in D'}\pair{z}{p}+\|\varphi-p\|
\ge\pair{e}{p}+\pair{e}{\varphi-p}=1.
\]
The choice $p=0$ attains $1$, proving the second equality in
Eq.~\eqref{eq:cor-PCT-gap}. Choose graph points over $0$ and $3e/2$
and assign each weight $1/2$. Their mean primal coordinate is $3e/4$,
which lies strictly inside the unit ball. Its value under $\varphi$
is $3/4$, exceeding the supremum on the left of
Eq.~\eqref{eq:cor-PCT-gap}.

For $d'=r_0e$ one has $x_2+\mathcal S d'=0$.
Eq.~\eqref{eq:cor-transformed-strip} and $\|\varphi\|=1$
give the distance lower bound. The definition of the Fitzpatrick
function and the full bijection of graph rows give
\[
\Phi_{\widetilde M}(d',0)
=\sup_{(x,a)\in\gra M}\pair{-x}{a}
=\Phi_M(0,0)=\sigma.
\]
This proves Eq.~\eqref{eq:cor-finite-projection-hole}.
\end{proof}
The two quantities in Eq.~\eqref{eq:cor-PCT-gap} are computed directly.
Equation~\eqref{eq:cor-finite-projection-hole} also gives a point
$(3e/4,0)$ where the Fitzpatrick function is finite although $3e/4$
lies outside the norm closure of the domain.

\begin{corollary}[Convex combinations of graph points and incompatible polar points]
\label{cor:incompatible}
For each original rank-one pair $(M,P_k)$, put
$G=\gra(M+P_k)$. We construct two points of $\gra M$ and one point
of $\gra P_k$ with the same mean primal coordinate. Their pairings produce
a point $q'$ and an explicit $\eta>0$ such that
\begin{equation}
q_0=(0,0)\in G^\mu,\qquad
[q',g]\ge\eta\quad(g\in G),\qquad
[q_0,q']=-\eta.
\label{eq:cor-incompatible}
\end{equation}
Here $[(x,a),(y,b)]=\pair{x-y}{a-b}$.
Thus $G^\mu$ is not monotone, and $G$ has at least two distinct,
incomparable maximally monotone extensions.
\end{corollary}
\begin{proof}
Set $t=65/64$, $s=1/64$, and
\[
v=\sum_{n\ge1}\omega_s(n)e_{n,1},\qquad
W=\|\omega_s\|_2^2,\qquad E=t^2-1=\frac{129}{4096}.
\]
The vector $v$ is finitely supported. The single first tail coordinate
already gives
\[
W\ge\left(\frac{63}{256}\right)^2
  =\frac{3969}{65536}>\frac{2064}{65536}=E.
\]
Take the two actual labels
\[
a^+=-t e_{0,1}+(t+1)e_{0,3}+v,
\qquad
a^-=t e_{0,1}-(t+1)e_{0,3}+v.
\]
They realize parameters $t$ and $-t$. On $c_0(I)$ put
$x^\pm=La^\pm$; their equal-weight mean is
$w=Lv=-\sum_n\omega_s(n)e_{n,1}$, and their dual mean is $v$.
One has $\pair{w}{g}=0$, so $(w,0)$ is an actual $P$-row.
Moreover,
\[
\pair{x^\pm}{a^\pm}=E-W,\qquad
\pair{w}{v}=-W.
\]
The correction term for the mean of the two points of $\gra A$ is
\[
\frac14\pair{x^+-x^-}{a^+-a^-}
=\frac12\pair{x^+}{a^+}+\frac12\pair{x^-}{a^-}
  -\pair{w}{v}=E.
\]
The corresponding correction for the single point of $\gra P$ is zero,
and $\pair{-w}{v}=W>E$.

On standard $\ell^1$, choose actual lifts $u^\pm$ with
$Qu^\pm=x^\pm$, and use rows $(u^\pm,Q^*a^\pm)$.
Their mean $\bar u=(u^++u^-)/2$ has
$\pair{\bar u}{f}=0$, so $(\bar u,0)$ is an actual $B$-row.
All the preceding pairing equalities are preserved by
$\pair{u}{Q^*a}=\pair{Qu}{a}$. No particular choice of lift,
linear right inverse, or change of norm is needed.

Write $r=(w,v)$ in the first space and
$r=(\bar u,Q^*v)$ in the second. Averaging monotonicity with
the specified graph points proves
\[
[r,g]\ge-E\quad(g\in G),\qquad [q_0,r]=-W.
\]
Indeed, for a sum row $(x,a+b)$, the averaged $M$ brackets
give $\pair{r_X-x}{r_{X^*}-a}\ge-E$, and the single
partner bracket gives $\pair{r_X-x}{-b}\ge0$.
The core relatedness $q_0\in G^\mu$ can also be read directly
from $\pair{x}{a+P_kx}=2|\pair{x}{k}|^2-1-W(|\pair{x}{k}|-1)>0$.
For
\[
\theta=\frac{W-E}{2W},\qquad q'=\theta r,
\qquad \eta=\frac{(W-E)^2}{4W}>0,
\]
the three-point identity yields
\begin{align*}
[q',g]
&=(1-\theta)[q_0,g]+\theta[r,g]
                 -\theta(1-\theta)[q_0,r]\\
&\ge\theta\bigl((1-\theta)W-E\bigr)=\eta.
\end{align*}
Also $[q_0,q']=\theta^2[r,q_0]=-\eta$.
Thus both points belong to $G^\mu\setminus G$ and no monotone
extension of $G$ can contain both. Each monotone set
$G\cup\{q_0\}$ and $G\cup\{q'\}$ has a maximally monotone
extension: a chain of monotone supersets has a monotone union,
so Zorn's lemma applies. The two extensions differ, and neither
can properly contain the other by maximality.
\end{proof}
The graph points and their pairings are given above; the maximal
extensions in the final sentence are existence conclusions, not
two additional explicit operator formulas. The averaged pair $r$
is not asserted to be a factor row. This distinction is essential:
the finite averaging identity produces the new polar point through
the proved bracket calculation, not by identifying a barycentre
with an actual graph row.
